\subsection*{トピックと参考文献の対応}

\addcontentsline{toc}{subsection}{4 トピックと参考文献の対応}

この表は、SWEBOK第14章が示す「トピック対参考資料の行列」を学習用に簡略化したものである。詳しく学ぶときは、各トピックに対応する文献を読むとよい。

\begin{center}
\small
\par\smallskip\noindent\refstepcounter{table}\label{tab:ch14-05}表 \thetable: トピックと参考文献の対応の整理\par\smallskip
表 \thetable は、トピックと参考文献の対応の整理を比較・確認しやすい形で整理したものである。\par\smallskip
\begin{longtable}{p{0.40\linewidth}p{0.55\linewidth}}
\toprule
トピック & 主な参考資料 \\
\midrule
\endfirsthead
\toprule
トピック & 主な参考資料 \\
\midrule
\endhead
1. 専門性 & Bott et al.; Sommerville; McConnell; Tockey \\
\midrule
1.1 認定・資格証明・資格・免許 & Bott et al.; Sommerville; Washington Accord; ISO/IEC 24773-1; ISO/IEC 24773-4 \\
\midrule
1.2 倫理規程と専門職行動規範 & Bott et al.; Voland; Sommerville; McConnell; ACM Code; IEEE Code; IFIP Code; Tockey \\
\midrule
1.3 専門職団体 & Bott et al.; Sommerville; McConnell \\
\midrule
1.4 ソフトウェア工学標準 & Bott et al.; Appendix B; Sommerville; McConnell \\
\midrule
1.5 経済的影響 & Voland; Sommerville; Tockey \\
\midrule
1.6 雇用契約 & Bott et al.; ACM Code; IEEE Code; IFIP Code \\
\midrule
1.7 法的論点 & Bott et al.; Appendix B; Voland; Sommerville \\
\midrule
1.8 文書化 & Bott et al.; Voland; Sommerville; McConnell \\
\midrule
1.9 トレードオフ分析 & Voland; Sommerville; Tockey \\
\midrule
2. 集団力学と心理学 & Voland; Sommerville; McConnell; Fairley \\
\midrule
2.1 チーム・集団で働く力学 & Voland; Fairley \\
\midrule
2.2 個人の認知 & Voland; McConnell \\
\midrule
2.3 問題の複雑さへの対応 & Voland; Sommerville; McConnell \\
\midrule
2.4 利害関係者との相互作用 & Sommerville \\
\midrule
2.5 不確実性と曖昧さ & Sommerville; Fairley \\
\midrule
2.6 公平性・多様性・包摂性 & Sommerville; Fairley \\
\midrule
3. コミュニケーション技能 & Voland; Sommerville; McConnell; Fairley \\
\midrule
3.1 読む・理解する・要約する & Sommerville; McConnell \\
\midrule
3.2 書く & Voland; Sommerville \\
\midrule
3.3 チームと集団のコミュニケーション & Voland; Sommerville; McConnell; Fairley \\
\midrule
3.4 プレゼンテーション技能 & Voland; Sommerville; Fairley \\
\midrule
\bottomrule
\end{longtable}
\normalsize
\end{center}
